\documentclass[english]{article}
\usepackage[T1]{fontenc}
\usepackage[latin9]{inputenc}
\usepackage[a4paper]{geometry}
\geometry{verbose,tmargin=3cm,bmargin=2.5cm,lmargin=2.5cm,rmargin=2.5cm}
\usepackage{color}
\definecolor{note_fontcolor}{rgb}{0.800781, 0.800781, 0.800781}
\usepackage{array}
\usepackage{verbatim}
\usepackage{rotfloat}
\usepackage{url}
\usepackage{amstext}
\usepackage{amsthm}
\usepackage{graphicx}
\usepackage{setspace}
\doublespacing

\makeatletter



\providecommand{\tabularnewline}{\\}

\newenvironment{lyxgreyedout}
  {\textcolor{note_fontcolor}\bgroup\ignorespaces}
  {\ignorespacesafterend\egroup}


\theoremstyle{plain}
\newtheorem{thm}{\protect\theoremname}
\ifx\proof\undefined
\newenvironment{proof}[1][\protect\proofname]{\par
\normalfont\topsep6\p@\@plus6\p@\relax
\trivlist
\itemindent\parindent
\item[\hskip\labelsep\scshape #1]\ignorespaces
}{%
\endtrivlist\@endpefalse
}
\providecommand{\proofname}{Proof}
\fi

%%%%%%%%%%%%%%%%%%%%%%%%%%%%%% User specified LaTeX commands.
\usepackage{amsmath,epsfig,natbib,pstricks,pst-text,pst-node,pst-plot,amssymb,wrapfig,threeparttable,rotating}

\makeatother

\usepackage{babel}
\providecommand{\theoremname}{Theorem}

\begin{document}

\title{Non-linear association joint models for clinical studies}

\author{Sable Walters-Somers 21693073\vspace{0.5cm}\\
WST795 Research Report\vspace{2cm}\\
Submitted in partial fulfillment of the degree BSc(Hons) Mathematical
Statistics \\
\vspace{0.5cm}\\
Supervisor(s): J. van Niekerk\vspace{2cm}\\
Department of Statistics, University of Pretoria\vspace{0.5cm}\\
\includegraphics[width=5cm]{UPlogohighres.jpg}\vspace{0.5cm}\\
}

%\date{}

\maketitle
\newpage{}
\begin{abstract}
Short summary of the research proposed. This should be a two to three
paragraphs long and should fully describe the content and contributions
of the research report.
\\

Joint models in a clinical context aim to simultaneously model a patients longitudinal biomarker data and their Time-To-Event (TTE) survival data. For the sake of statistical, mathematical, theoretical, and practical simplicity, a linear linkage function between the two processes is assumed. However, this may easily fail to capture the complex and often non-linear processes and mechanisms that characterize pharmacokinetics and pharmacodynamics.
\\

This study aims to investigate the the effects of using a non-linear association function in a joint model within a clinical context; advantages, disadvantages, and general efficacy will be investigated and discussed.

\newpage{}
\end{abstract}

\section*{Declaration}

I, \emph{Sable Walters-Somers}, declare that this research report, submitted in
partial fulfillment of the degree \emph{BSc(Hons) Mathematical Statistics}, at
the University of Pretoria, is my own work and has not been previously
submitted at this or any other tertiary institution. 

\vspace{1cm}

\_\_\_\_\_\_\_\_\_\_\_\_\_\_\_\_\_\_\_\_\_\_\_\_\_\_\_\_\_

\emph{Sable Walters-Somers}

\vspace{1cm}

\_\_\_\_\_\_\_\_\_\_\_\_\_\_\_\_\_\_\_\_\_\_\_\_\_\_\_\_\_

\emph{Janet van Niekerk}

\vspace{1cm}

\_\_\_\_\_\_\_\_\_\_\_\_\_\_\_\_\_\_\_\_\_\_\_\_\_\_\_\_\_

Date

\begin{comment}
Edit all italics with your details and sign on the solid lines once
the final document is printed and submitted. 
\end{comment}



\newpage{}

\tableofcontents{}

\listoffigures

\listoftables

\newpage{}

\section{Introduction}

YOUR 10 PAGE COUNT STARTS HERE

The introduction should provide a detailed description of the topic
and the aim of the research report. In addition the literature review
should intertwine with this. It is important to always reference where
needed. All work from somewhere else requires a reference \cite{Anguelov2010}.


QUESTIONS/POINTS TO COVER:

-

-

-

-

-

Inline equation: $x=y$

Display equation: \[x=y\]

Numbered display equation: 
\begin{equation}
x=y\label{eq:afe,fjgd}
\end{equation}

Numbered equation: 
\begin{equation}
x=y+1
\end{equation}

Equation array: 

\begin{eqnarray*}
x & = & y+2+3\\
 & = & y+5.
\end{eqnarray*}


\section{Background Theory}

In Figure \ref{fig:UPlogo} and in Figure \ref{fig:Caption} we see
the UP logo. %

\begin{figure}[h]
\centering
\includegraphics[width=6cm]{UPlogohighres.jpg}
\caption{A graphic of the UP logo\label{fig:UPlogo}}
\end{figure}

- What should be in a caption? 

- How should I get my graphics?

\begin{figure}
\includegraphics[width=6cm]{UPlogohighres.jpg}
\caption{Caption \label{fig:Caption}}
\end{figure}


 The theory of the topic should be thoroughly discussed in this section.
The student must show their proficiency on the topic as well as additional
insight. This section may be separated into a few sections as \textit{necessary}.
Note that the theory of the topic is the main contribution and the
report should illustrate this. In Theorem \ref{thm:first} we discuss.....
$f(x)=\left\{ \begin{array}{cc}
1 & \textrm{\textrm{if }}x<2\\
2 & \textrm{if }x\ge2
\end{array}\right.$
\begin{thm}
\label{thm:fbkndfblksnfbkn}fbkndfblksnfbkn 

jbj fbfjk b bfjf $f(x)=\left\{ \begin{array}{cc}
1 & \textrm{\textrm{if }}x<2\\
2 & \textrm{if }x\ge2
\end{array}\right.$

\end{thm}
\begin{proof}
blksnfblkdfsnb
\end{proof}
m,fgkrjtbnel
\begin{thm}
\label{thm:first} aurbyveurybverybyruwei 
\end{thm}
\begin{proof}
vnitunrivuptnriuwvtnoiutnreuvnt
\end{proof}



\section{Application}

\begin{sidewaystable}
\begin{centering}
\begin{tabular}{|>{\centering}p{3cm}|c|c|c|>{\centering}p{3cm}|}
\multicolumn{1}{>{\centering}p{3cm}|}{} &  &  &  & \tabularnewline
klvunl ruvynryrinfsfgng hdfhfgjf gjfhjghjhk jhkghjkh dfhdj ghkjfgj
djdj dghjhhjghj & aevirue$x=y$n oiurlevniq & ermwovmi;q & aevpimqw & ovmi;\tabularnewline
 &  &  &  & \tabularnewline
 &  &  &  & \tabularnewline
 &  &  &  & \tabularnewline
\end{tabular}
\par\end{centering}
\caption{rvybruvybrv \label{tab:test}}

\end{sidewaystable}

\begin{table}
\begin{centering}
\begin{tabular}{|>{\centering}p{4cm}|>{\centering}p{4cm}|>{\centering}p{4cm}|}
\hline 
\multicolumn{1}{|>{\centering}p{4cm}||}{} & \multicolumn{1}{>{\centering}p{4cm}||}{Column 1} & Column 2\tabularnewline
\hline 
\hline 
Row 1 &  & \tabularnewline
\hline 
Row 2 &  & \tabularnewline
\hline 
\end{tabular}
\par\end{centering}
\caption{Example \label{tab:Example}}
\end{table}
 The application should be presented in this section. Code should
be included in an appendix as well as additional output if needed. 

\section{Conclusion}

The conclusion should summarise what was done in the research report.
It should also provide shortfalls of the research and recommendations
on what could be investigated in future. This section should be an
honest summary of the research.

YOUR 10 PAGE COUNT ENDS HERE

\newpage

\bibliographystyle{plain}
\bibliography{BibliographyDatabase}

\newpage

\section*{Appendix\addcontentsline{toc}{section}{Appendix}}

Include any additional code, output or data here. 
\end{document}
